\documentclass[11pt,a4paper]{article}
\usepackage[margin=2.4cm]{geometry}
\usepackage{booktabs}
\usepackage{array}
\usepackage{xcolor}
\usepackage{longtable}
\usepackage{enumitem}
\usepackage{xurl}
\usepackage[hidelinks,colorlinks=true,linkcolor=black,urlcolor=blue!60!black,citecolor=black]{hyperref}
\usepackage{titlesec}
\usepackage{fancyhdr}

\definecolor{ink}{HTML}{0B0B0B}
\definecolor{muted}{HTML}{52514E}
\definecolor{rule}{HTML}{C9C7C2}

\titleformat{\section}{\large\bfseries}{\thesection}{0.7em}{}
\titleformat{\subsection}{\normalsize\bfseries}{\thesubsection}{0.6em}{}
\setlist[itemize]{leftmargin=1.2em,itemsep=2pt,topsep=3pt}
\setlist[enumerate]{leftmargin=1.4em,itemsep=2pt,topsep=3pt}
\renewcommand{\arraystretch}{1.25}

\pagestyle{fancy}\fancyhf{}
\fancyhead[L]{\footnotesize\color{muted}Weekend Margin Clearinghouse — Going to Market}
\fancyhead[R]{\footnotesize\color{muted}Team 9}
\fancyfoot[C]{\footnotesize\color{muted}\thepage}
\renewcommand{\headrulewidth}{0.4pt}

\newcommand{\note}[1]{{\small\color{muted}#1}}

\begin{document}

\begin{center}
{\LARGE\bfseries Going to Market}\\[4pt]
{\large An on-chain clearinghouse for perpetual swaps margined with tokenized T-bills}\\[10pt]
{\normalsize FRE-GY 6931 Blockchain \& Cryptocurrency, Fall 2026 — Team 9}\\[2pt]
{\small Sai Kishore (contracts, on-chain) \quad·\quad Swastik Mishra (analysis, product)}\\[2pt]
{\small Prof.\ Richard Radnay \quad·\quad TA Hina Tomar \quad·\quad \today}
\end{center}
\vspace{4pt}
\hrule
\vspace{14pt}

\section*{Executive summary}
\addcontentsline{toc}{section}{Executive summary}

Between 2024 and 2026 the market moved decisively toward tokenized Treasuries as derivatives
collateral. Circle launched a 24/7 BUIDL-to-USDC contract in April 2024 \cite{circle}. Binance
accepted Circle's USYC as off-exchange collateral in July 2025 and BlackRock's BUIDL that November
\cite{binance,protofire}. Deribit and Crypto.com followed \cite{deribit}. In September 2026, CFTC
staff confirmed that swap dealers may post tokenized money market fund shares as margin
\cite{protofire}.

Every one of those moves validates the premise this project was built on: institutions want to post
interest-bearing collateral against derivatives, and the settlement calendar is what stands in the
way.

\textbf{What none of them does is unify the two halves.} Tokenized funds are issued without a
clearing layer, and clearing venues operate without owning the collateral. BlackRock's fund does not
know it is being used as margin; Deribit does not know what is inside it. The gap between them is
bridged by intermediaries --- Ceffu, triparty banks, Circle's discretionary USDC reserve --- and each
of those reintroduces precisely the counterparty risk and operating-hours constraint that
tokenization was meant to remove.

\textbf{We build both halves as one system.} Our clearinghouse natively understands its collateral's
redemption calendar. When the issuer is unreachable, \texttt{settle()} routes through a permissionless
vault that buys the collateral at a discount and redeems it when the issuer reopens. No custodian, no
discretionary counterparty, no approval queue. The asset and the clearing layer are the same system
and speak the same language, which is what makes automatic settlement possible at three in the
morning on a Sunday.

That is the novelty, and it determines the product. This is \textbf{non-custodial settlement liquidity
for on-chain derivatives} --- built for venues where no central balance sheet exists to absorb a
timing mismatch, which is exactly where the incumbent model cannot reach without becoming an
intermediary again.

Three quantitative conclusions follow, each developed below:

\begin{enumerate}
  \item \textbf{Break-even is \$134m of hedged notional}, supported by a \$26m vault — that vault
  being about 0.2\% of the \$15bn tokenized treasury market.
  \item \textbf{The deployed 10bp fee is not commercially viable.} The vault earns 1.21\% a year
  against 4.01\% on T-bills. A real product needs 28--44bp, and the band where both sides agree is
  narrow: 28--42bp.
  \item \textbf{The design is economically impossible on Ethereum mainnet.} At 100 open trades the
  keeper's gas alone costs \$1.90m a year, about sixty times the revenue those trades earn. On an L2 the same
  workload costs \$2,078. This single constraint determines where the product can exist.
\end{enumerate}

\section{Market context}

Tokenized U.S.\ Treasury products held roughly \$14.5--16bn in October 2026, growing about 125\% year
over year \cite{allium,cryptonomist,utoday}. Allium counts 91 distinct assets from 36 issuers
\cite{allium}. Circle's USYC is now reported as the largest single fund at approximately
\$2.6--3bn, having overtaken BlackRock's BUIDL \cite{coincentral}.

\note{Trackers disagree because on-chain token market capitalisation differs from fund AUM. We quote
a range rather than a point estimate and cite three sources.}

Two regulatory and market developments matter for us. First, derivatives venues began accepting these
instruments as margin collateral: Binance added USYC in July 2025 and BUIDL in November 2025, holding
units in off-exchange custody at Ceffu or a triparty bank \cite{binance,protofire}; Deribit and
Crypto.com followed, with Deribit reported to apply a 2\% haircut \cite{deribit,protofire}. Second,
CFTC staff guidance issued 24 September 2026 is reported to confirm that swap dealers may post
tokenized money market fund shares as margin \cite{protofire}.

\note{We could not confirm the CFTC item against a primary CFTC document. See Appendix B.}

\section{Competitive landscape}

\subsection{Three incumbent models}

\textbf{Venue collateral (Binance, Deribit, Crypto.com).} The fund posts tokenized shares, accepts a
haircut, and receives margin credit. When it loses money on a Saturday, the venue debits its
collateral balance as an internal ledger entry. \emph{No cash moves.} The venue absorbs the timing
mismatch on its own balance sheet and settles later.

\textbf{Issuer off-ramp (Circle for BUIDL).} A single counterparty holds a USDC reserve and purchases
shares around the clock \cite{circle}. It is structured as a secondary-market transfer rather than a
fund redemption, so accrued interest still pays at the next distribution cycle \cite{techflow}.
Reported reserve capacity is of the order of \$80--100m \cite{techflow,bitcoincom}.

\textbf{Primary redemption (Securitize).} The conventional path: daily, business days, with a 3:00pm
cutoff \cite{techflow}. This is the constraint our project was designed around, and it remains real.

\subsection{Comparison}

\begin{center}
\small
\begin{tabular}{@{}p{4.3cm}p{2.9cm}p{2.9cm}p{3.5cm}@{}}
\toprule
& \textbf{Circle off-ramp} & \textbf{Venue collateral} & \textbf{Weekend Vault (ours)} \\
\midrule
Liquidity provider & Circle & The venue & Any depositor \\
Permission required & Securitize allowlist & Venue onboarding & None \\
Discretionary & Yes & Yes & No — contract logic \\
Capacity & \textasciitilde\$80--100m, one party & Venue balance sheet & Scales with deposits \\
Serves on-chain perps & No & No & \textbf{Yes} \\
Multi-issuer & BUIDL only & Per venue & Any ERC-20 with a NAV \\
Fires automatically at margin call & No & N/A (ledger entry) & \textbf{Yes} \\
Counterparty risk & Circle & Venue + custodian & Smart contract \\
Operating history & \textasciitilde2.5 years & 1--2 years & \textbf{None} \\
\bottomrule
\end{tabular}
\end{center}

\subsection{The structural gap}

Incumbent solutions work because a trusted party absorbs the mismatch, and for users who can reach
that party they work well. The limitation is structural rather than one of execution: the model
requires an intermediary, so it cannot serve any venue that does not have one.

Derivatives that clear on-chain are exactly that case. Our clearinghouse is bilateral with no central
balance sheet, so a loss on a Saturday cannot be settled by moving a number in a ledger --- the cash
has to move. That constraint is what produces the vault, and it is the position incumbents cannot
occupy while remaining intermediaries.

The deeper division is between layers. Issuance and clearing are built by different firms, on
different infrastructure, under different regulators, and are stitched together after the fact by
custodians and bilateral agreements. We treat them as one problem: the clearinghouse reads the
collateral's liquidity state directly and routes settlement accordingly, with the vault as the
counterparty of last resort. The course's requirement that every project mint and use its own token
pushed us toward this design, and it turned out to be the more interesting half of the work.

\section{Users}

We identify four segments, ordered by how poorly incumbents serve them.

\begin{enumerate}
  \item \textbf{On-chain perpetual traders} (Hyperliquid, dYdX, GMX and similar). No central balance
  sheet exists to absorb a weekend mismatch, so settlement requires real cash movement. This is the
  primary market.
  \item \textbf{Bilateral OTC derivatives between funds.} No clearer, same structural problem.
  \item \textbf{Funds below BUIDL's access threshold.} BUIDL is offered under Section 3(c)(7) of the
  Investment Company Act to qualified purchasers, with a reported \$5m minimum
  \cite{ecokyc,coinpaprika}. Everyone beneath that line cannot use Circle's off-ramp at all.
  \item \textbf{Holders of issuers nobody has integrated.} Ninety-one tokenized treasury products
  exist \cite{allium}; the Circle off-ramp serves one.
\end{enumerate}

\textbf{What they get that they do not have today:} the ability to meet a variation margin call in
cash, at any hour, without pre-funding an idle cash buffer and without asking anyone's permission.
Section~\ref{sec:business} quantifies the benefit at up to 42bp against the realistic alternative.

\section{Trust}

\subsection{What a stranger can verify}

All 16 deployed contracts are verified on Etherscan, confirmed through the Etherscan API, and the
verified source matches the files in this repository. The test suite runs with one command and
contains 37 tests, 25 of which expect a revert, plus three seeded randomised invariant runs checking that no tUSDC or
sBILL is created or destroyed. OpenZeppelin is used throughout for the token, access control,
\texttt{ReentrancyGuard} and \texttt{Pausable}.

The code has been scanned with \textbf{Slither 0.11.6} against the deployment compiler. It returned
26 findings on our own files. We fixed four — an uninitialised local in \texttt{liquidate}, two
shadowed \texttt{owner} parameters, and an unindexed event address — and \texttt{docs/slither.md}
explains each of the remaining 22 as a false positive or an accepted design choice, finding by
finding. Independently, 28 values read from the live Sepolia contracts reconcile to the wei against
the financial model.

\textbf{sBILL is a mock.} It accrues a fixed 4\% and pays interest from a pre-funded reserve, so it demonstrates the integration rather than carrying real T-bill exposure. Section~\ref{sec:risks} sets out what a production issuer would require.

\subsection{Admin powers, stated exactly}

The owner may pause \emph{new} deposits and \emph{new} trades, and may post a funding rate within a
cap fixed at deployment (0.1\% per hour). The owner \textbf{cannot} interfere with settlement,
liquidation, trade closing or withdrawals, all of which continue while paused.

\texttt{MarketCalendar} has no owner and no settings. Nobody, including us, can change when the
weekend is. This was a deliberate choice against a specific attack: an admin able to repoint the
calendar at a contract reporting ``always closed'' could force every settlement through the vault at
a discount, brick fund redemptions, and raise the collateral haircut toward liquidation thresholds —
three extraction routes from one key.

\textbf{We disclose that the demonstration deployment is different.} It uses a separately labelled
\texttt{DemoCalendar} whose weekend is switched by the deployer, so the mechanism can be shown on any
day. The production design is the admin-free calendar.

\subsection{What the interface does to stop costly mistakes}

A clearinghouse that is correct on-chain and opaque in the browser is not trustworthy. The front end
carries four specific protections, each of which can be inspected in the repository.

\begin{itemize}
  \item \textbf{Readable error messages.} \texttt{utils/weekend/errors.ts}
  maps 27 custom errors — ours and OpenZeppelin's — to sentences that say what happened and what to
  do, plus fallbacks for a rejected signature and insufficient gas. It is wired into the shared
  transaction path at \texttt{utils/scaffold-eth/contract.ts:345}, so it covers every call the
  application can make rather than a hand-picked few. A user who posts too large a trade reads
  ``Your margin is too low for this. Deposit more, or use a smaller size,'' not
  \texttt{InsufficientMargin(1, 2)}.
  \item \textbf{Clear confirmations on irreversible actions.} The shared button component takes a
  \texttt{confirmText}; the first click only asks, so the action needs two clicks. It guards the
  four actions that cannot be undone: opening a trade, closing one, liquidating a counterparty,
  and pausing the system. Closing reads
  ``This settles at the current price and ends the trade for both sides.''
  \item \textbf{Impossible actions are disabled with a stated reason, not left to fail on-chain.}
  There are 24 \texttt{blockedBecause} guards across the five application pages and the shared
  components. A button that cannot work is greyed out and says why — no wallet connected, the price
  is stale, the amount is empty, owner-only — so the user never pays gas to discover a revert.
  \item \textbf{The liquidation price is shown before the trade is opened,} not after. The trade
  page computes ``You'd be liquidated around \$X'' from the same formula the contract uses, for both
  sides, as the size and margin inputs change. When the Chainlink price is stale or unset, a warning
  banner states that new trades and settlement are waiting.
  \item \textbf{Slippage limits are unnecessary here, by construction.} There is no pool and no
  order book to walk, so there is no price impact to bound. A trade opens at the Chainlink price both
  sides can read before committing, and a weekend sale to the vault executes at NAV minus a discount
  fixed in contract storage. The user's fill is deterministic and knowable in advance, which is a
  stronger guarantee than a slippage tolerance. The honest counterpart is that the price is only as
  good as the oracle, which is why the two-hour staleness check and the banner above exist.
\end{itemize}

\textbf{What we have not done.} Every flow was exercised on a local chain with a test wallet, and on
Sepolia the pages read the live contracts correctly, but no transaction has yet been sent from the
deployed site through MetaMask. We state this rather than imply coverage we do not have.

\subsection{Failure modes}

We state these rather than conceal them. All are safe failures — the system stops rather than
transacting on bad data — but each is a real limitation.

\begin{itemize}
  \item \textbf{Stale price feed.} Any Chainlink answer older than two hours is rejected. Settlement,
  liquidation and withdrawals from accounts with open trades revert until a fresh price arrives.
  \item \textbf{Vault exhausted.} If the vault lacks cash, \texttt{settle()} and \texttt{liquidate()}
  revert until cash arrives or the market reopens.
  \item \textbf{No insurance fund.} A sufficiently large gap move can leave a winner underpaid.
  \item \textbf{Holiday blindness.} The calendar is day-of-week only, so on a bond market holiday the
  contract believes the market is open. With a real tokenized fund, settlement would stick until the
  next business day. Section~\ref{sec:risks} and the report's design-decisions section develop this.
\end{itemize}

\section{Promotion: the first hundred users}

Generic community-building is not a plan here, because the binding constraint is not awareness.

\begin{itemize}
  \item \textbf{Issuer integration is step one, not marketing.} BUIDL's token contract blocks
  transfers to wallets not on a Securitize-administered allowlist \cite{ecokyc,coinpaprika}. The vault
  and clearinghouse must be approved holders before a single real user can exist. This is a business
  development conversation, and it gates everything else.
  \item \textbf{Launch where the allowlist is not.} Franklin Templeton's BENJI sets minimum initial
  investments by chain: \$20 on Stellar, \$100 on Aptos, \$1,000 on Polygon, \$20,000 on Avalanche and
  \$5,000,000 on Ethereum \cite{franklin}. Our Ethereum-based design sits in the most restrictive
  tier. A cheaper chain resolves the access barrier and the gas constraint simultaneously.
  \item \textbf{Seed the vault ourselves.} This is a two-sided market with a cold-start problem: no
  vault, no settlement guarantee; no traders, no fee income. The first \$1m of vault capital has to
  come from the operator.
  \item \textbf{Integrate rather than acquire.} The realistic route to users is as a settlement module
  inside an existing on-chain perpetuals venue, not as a destination users navigate to.
  \item \textbf{Incentives, pointed at the short side of the market.} This is a two-sided product and
  only one side is scarce. Hedging funds arrive on their own, because the alternative is an idle cash
  buffer costing them up to 42bp. Vault depositors do not: they are asked to accept T-bills plus
  roughly 40bp for contract risk, which is a thin offer from an operator with no history. So the
  incentive budget goes to depositors, not traders: the protocol waives its 25\% share of the
  discount for the first year, which lifts the depositor's return from 4.41\% to 5.76\% and costs us
  revenue we would not have earned anyway at sub-scale volume. Trading incentives would be spent on
  the side that does not need them.
\end{itemize}

\subsection{The token, and why it is not a campaign}

Our minted token is \textbf{sBILL}, the ERC-20 tokenized T-bill share. It is the collateral the
system exists to clear, not an instrument of promotion: it has a NAV, it is redeemed at that NAV, and
nothing about it is speculative by construction.

We deliberately do not issue a governance or reward token, and the reason is specific to this
product rather than a general objection. The entire value proposition is that a fund can post
collateral and be certain of what it will realise. A token that subsidises deposits makes the
vault's advertised return partly a function of that token's price, which reintroduces exactly the
volatility the customer came here to remove. A depositor earning ``4.41\% plus an uncertain amount of
a governance token'' cannot price the product against T-bills, and pricing against T-bills is the
only comparison that matters here.

The promotional asset is therefore the \textbf{live reconciliation}, not a token launch: 28 numbers
read off the deployed contracts matching the model to the wei, and a weekend settlement that anyone
can replay on Etherscan. For an audience of risk officers, a verifiable number is worth more than an
airdrop.

\section{Business model}\label{sec:business}

\subsection{Where the revenue comes from}

The vault buys collateral at NAV minus a discount while the issuer is unreachable and redeems at full
NAV when it reopens. It earns the discount plus the interest the asset accrues while held — tokenized
treasuries accrue continuously, including weekends, even though they cannot be redeemed. Per \$1 of
vault capital, with the pool sized to cover 99\% of weekends — which leaves it 24.7\% utilised on
average — and a 4.01\% T-bill yield:

\begin{center}
\small
\begin{tabular}{@{}lrrrp{3.5cm}@{}}
\toprule
\textbf{Discount} & \textbf{Income} & \textbf{Carry} & \textbf{Total} & \\
\midrule
10bp (as deployed) & 1.29\% & 0.36\% & \textbf{1.65\%} & loses to T-bills \\
28bp & 3.60\% & 0.36\% & \textbf{3.96\%} & matches T-bills \\
42bp (our price) & 5.40\% & 0.36\% & \textbf{5.76\%} & before our share \\
44bp & 5.66\% & 0.36\% & \textbf{6.02\%} & T-bills $+$ 2\% \\
\bottomrule
\end{tabular}
\end{center}

\textbf{At the deployed 10bp the vault loses to its own alternative}, earning 1.65\% against 4.01\% on
T-bills. The cause is not the weekend — it is that three quarters of the vault's capital is idle on a
normal weekend. Our financial analysis establishes the per-use cost at only 6.6bp, which 10bp covers
comfortably. The shortfall is entirely the cost of waiting, and it is why the contract was deployed
with an owner-adjustable discount: the owner can raise it up to a hard limit of 1\%, only while the
market is open, without redeploying.

\subsection{The band where a deal exists}

Our analysis derives both sides of the price. The vault requires 28--44bp serving both directions from
one pool. The fund will pay up to 119bp against holding everything in cash, 42bp against the realistic
alternative of a cash buffer, and only 16bp against a fund disciplined enough to sell T-bills every
Friday.

\begin{center}
\begin{tabular}{@{}lr@{}}
\toprule
\textbf{Constraint} & \textbf{Discount (bp)} \\
\midrule
Vault must earn at least (both directions, match T-bills) & 28 \\
Vault must earn at least (both directions, T-bills + 2\%) & 44 \\
Fund will pay at most (vs.\ cash buffer) & 42 \\
Fund will pay at most (vs.\ weekly T-bill sales) & 16 \\
\bottomrule
\end{tabular}
\end{center}

\textbf{A workable price exists between 28 and 42bp, and only against the cash-buffer alternative.}
Against a weekly-sale strategy there is no overlap at all. This is a narrow business, and we report it
as such.

\subsection{What we charge, and who it has to satisfy}

Three parties must all be satisfied by one number, and the window is narrow.

\begin{center}
\begin{tabular}{@{}p{8.6cm}r@{}}
\toprule
\textbf{Constraint} & \textbf{Discount (bp)} \\
\midrule
Vault must earn at least this to match T-bills (both directions, one pool) & 28 \\
Vault must earn at least this for T-bills $+$ 2\% & 44 \\
\textbf{Fund will pay at most this, against a cash buffer} & \textbf{42} \\
Fund will pay at most this, against weekly T-bill sales & 16 \\
\bottomrule
\end{tabular}
\end{center}

\textbf{We price at 42bp, the top of the window, and the protocol takes 25\% of it.} That leaves
depositors 31.5bp, which the model turns into \textbf{4.41\% a year against 4.01\% on T-bills}.

Two things follow, and we state both. First, the fund's 42bp ceiling sits \emph{below} the 44bp the
vault would need to pay depositors a 2\% risk premium — so at no price does this product clear a
conventional hurdle rate. It clears a T-bill-flat hurdle with about 40bp to spare, and that is the
honest description of the opportunity. Second, pricing at the customer's ceiling leaves no room to
absorb a rate move, which is why Section~\ref{sec:risks} treats rate dependence as a first-order
risk rather than a footnote.

\subsection{Break-even volume}

Costs are modelled explicitly rather than as a range.

\begin{center}
\begin{tabular}{@{}lrp{7.2cm}@{}}
\toprule
\textbf{Annual cost} & \textbf{Amount} & \textbf{Basis} \\
\midrule
Infrastructure and keepers & \$24{,}000 & Hosting, RPC and monitoring. Keeper gas itself is only about \$1{,}900 a year at 100 open trades on an L2 \\
Security audits & \$60{,}000 & Mid-range for a focused DeFi protocol; published 2026 guidance spans roughly \$15k--\$100k \cite{sherlock} \\
Legal and compliance & \$120{,}000 & Entity, counsel, and the issuer agreements Section~5 shows are mandatory \\
Team & \$300{,}000 & Two engineers \\
\midrule
\textbf{Total} & \textbf{\$504{,}000} & \\
\bottomrule
\end{tabular}
\end{center}

Revenue has two legs. The vault leg is the protocol's share of the discount actually earned, which
depends on how much collateral is sold through the vault — on average 2.50\% of a hedge's notional
per weekend. The clearing leg is 1bp per side, charged each time a position is opened or rolled,
assumed twelve times a year. That fee sits well under the 4.5--5bp taker fees at Hyperliquid and
dYdX, which is appropriate: we are a settlement layer beneath a venue, not the venue.

\begin{center}
\begin{tabular}{@{}lr@{}}
\toprule
\textbf{Per \$1m of hedged notional, per year} & \\
\midrule
Vault fees (25\% of a 42bp discount) & \$1{,}366 \\
Clearing fees (1bp $\times$ 2 sides $\times$ 12 turns) & \$2{,}400 \\
\midrule
\textbf{Total} & \textbf{\$3{,}766} \\
\bottomrule
\end{tabular}
\end{center}

\textbf{Break-even is \$134m of hedged notional.} At a \$200m target the protocol earns \$753{,}277
against \$504{,}000 of costs, and the vault behind it needs \$39.4m of cash — 19.7\% of notional,
the level that covered 99\% of historical weekends when longs and shorts are served from one pool.
Pooling the two sides rather than running separate books saves 29\% of the capital, which is the
single largest efficiency in the design.

Note which leg carries the business: \textbf{clearing fees are 64\% of revenue}. The vault is what
makes the product possible, but a flat per-trade fee is what pays for it. A reader who disbelieves
our discount analysis entirely still finds a business at \$210m of notional on clearing fees alone.

\subsection{The gas constraint}

Our keeper settles each open trade hourly. This cost is absent from our financial analysis and it is
decisive. Using the measured cost of a settlement paid in tUSDC, 72,188 gas (\texttt{docs/gas.md},
median from \texttt{yarn test}), at 10 gwei and \$3,000 ETH:

\begin{center}
\begin{tabular}{@{}lrr@{}}
\toprule
\textbf{Open trades} & \textbf{Ethereum mainnet} & \textbf{L2 (\textasciitilde1000\(\times\) cheaper)} \\
\midrule
2 & \$37,942 & \$38 \\
100 & \textbf{\$1,897,101} & \$1,897 \\
1,000 & \$18,971,006 & \$18,971 \\
\bottomrule
\end{tabular}
\end{center}

At 100 open trades, mainnet gas alone costs \$1.90m a year. Those trades are perhaps \$8m of
notional at the demo's size, which earns about \$30{,}000 of revenue — the gas is sixty times the
income. \textbf{The product cannot exist on Ethereum mainnet.} On an L2 the identical workload costs
\$1,897, which is immaterial, and it is the reason the \$24{,}000 infrastructure line in
Section~\ref{sec:business} is dominated by hosting rather than by gas.

The same arithmetic decides whether liquidation works. The reward is 50bp of notional and
\texttt{liquidate} costs 170,636 gas, so on mainnet a liquidator breaks even at about \$1{,}024 of
notional and below that no rational keeper acts. On an L2 the threshold is about \$1, so the
incentive holds at any size a user would trade.

This is the single most consequential finding in our go-to-market work. It does not change the
contracts, but it determines where they can be deployed, and it compounds with the BENJI access tiers
in Section~5: the cheaper chain is both the regulatory entry point and the only viable gas regime.

\section{Risks}\label{sec:risks}

\begin{enumerate}
  \item \textbf{Incumbents extend downward.} If Circle widens its off-ramp beyond BUIDL, or on-chain
  venues negotiate their own issuer arrangements, the gap we occupy closes. We consider this the most
  likely adverse outcome.
  \item \textbf{Holder restrictions.} Tokenized funds restrict who may hold them. BUIDL enforces a
  Securitize allowlist at the token contract \cite{ecokyc,coinpaprika}; BENJI shares are held in
  KYC'd registered accounts with chain-dependent minimums \cite{franklin}. Our vault and
  clearinghouse would both require approved-holder status, and vault depositors might require KYC.
  Our design limits the damage: sBILL moves only to the vault or back to the fund, never to a
  liquidator, so no unapproved transfer path exists.
  \item \textbf{Thin economics.} The viable band is 28--42bp and disappears entirely against a fund
  that sells T-bills weekly.
  \item \textbf{Interest rate dependence.} At 2021 rates of 0.05\%, a user of this product would have
  \emph{lost} approximately \$2,222 a year on a \$1m hedge. The product requires meaningful
  short-term rates.
  \item \textbf{Gas regime.} Mainnet is not viable, so the business inherits the security and
  liquidity assumptions of whichever L2 it deploys to.
  \item \textbf{Regulatory ambiguity.} The September 2026 CFTC staff guidance reportedly legitimises
  tokenized money market fund collateral for \emph{swap dealers} — that is, for the intermediated
  model we are positioned against \cite{protofire}.
  \item \textbf{A market downturn exhausts the vault before it threatens the fee.} Our regime
  testing puts numbers on this. Sizing the vault at the 95th percentile of weekend moves gives
  \$93{,}522 of capacity per \$1m of notional, and \textbf{four of the seven historical crises we
  tested exceed it} — the August 2024 yen-carry unwind needed 218\% of capacity, the June 2022 CPI
  shock 204\%, the SVB weekend 145\% and the May 2021 mining ban 137\%. When the vault empties,
  \texttt{settle()} reverts and settlement waits for Monday. This is the binding constraint on the
  business, and it is a capital question rather than a pricing one: the fee is comfortable in every
  one of those weekends, and the cash is not.
  \item \textbf{Exploit.} The contracts are unaudited. Slither reports 26 findings, of which we
  fixed four and explain 22; the single High-severity item is a false positive, where the tool reads
  the modulo in \texttt{MarketCalendar} as weak randomness when it is arithmetic on a timestamp. Four
  \texttt{reentrancy-no-eth} findings remain by choice: state is written after external calls, but
  every entry point is \texttt{nonReentrant}, the callees are our own contracts fixed at deployment,
  and the ERC-20s here have no transfer hooks. That reasoning is sound for our tokens and would need
  revisiting against a production issuer whose token may have hooks or a blocklist. An external audit
  is budgeted at \$60{,}000 in Section~\ref{sec:business} and has not been performed.

  Two further weaknesses we found ourselves rather than inherited. \texttt{WeekendVault} emits its
  redemption profit as an unsigned value, so a loss is reported as zero and a depositor reading the
  event log cannot distinguish breaking even from losing money; emitting a signed
  \texttt{int256} fixes it. And because vault inventory is held at cost, a depositor who times a
  deposit just before Monday's redemption collects a share of a profit they bore no risk for. Neither
  is an exploit in the sense of letting an attacker take funds, and both are real.
  \item \textbf{Operating history.} We have none. Every competitor named here has one to
  two years of production operation. Our evidence is a single live run: one weekend settlement
  through the vault, one liquidation against a mock price, and 28 reconciled values.
\end{enumerate}

\appendix
\section{Assumptions}

\begin{center}
\small
\begin{tabular}{@{}p{4.1cm}p{3.3cm}p{7.6cm}@{}}
\toprule
\textbf{Assumption} & \textbf{Value} & \textbf{Basis} \\
\midrule
Vault average utilisation & 24.7\% & Measured over 304 weekends, pool sized at the 99th percentile \\
Vault sizing & 19.7\% of notional & 99th percentile, pooled both directions \\
Discount charged & 42bp & The fund's ceiling against a cash buffer \\
Protocol share of discount & 25\% & Chosen so depositors still clear T-bills; the ceiling is 36\% \\
Clearing fee & 1bp per side, 12 turns & Below the 4.5--5bp taker fees at Hyperliquid and dYdX \cite{hlfees} \\
Operating cost & \$504k/yr & Itemised in \S\ref{sec:business}; team is \$300k of it \\
Settlement gas & 72,188 & Measured, \texttt{yarn test} median, \texttt{docs/gas.md} \\
Gas price / ETH & 10 gwei / \$3,000 & Illustrative \\
L2 cost factor & \textasciitilde1000\(\times\) cheaper & Order-of-magnitude \\
T-bill yield & 4.01\% & FRED DTB3, 28 Oct 2026 \\
Vault hurdle & T-bills, not T-bills $+$ 2\% & A 2\% premium needs 44bp, above the fund's 42bp ceiling \\
\bottomrule
\end{tabular}
\end{center}

\section{Source verification status}

Much of the competitive information is secondary reporting and sources conflict. We record
confidence explicitly rather than presenting all claims as equally established.

\begin{center}
\small
\begin{tabular}{@{}p{7.0cm}lp{4.3cm}@{}}
\toprule
\textbf{Claim} & \textbf{Confidence} & \textbf{Action} \\
\midrule
Circle BUIDL\(\rightarrow\)USDC contract exists, April 2024 & High & Primary press release \cite{circle} \\
BUIDL restricted to qualified purchasers, allowlisted & High & Multiple sources agree \\
BENJI chain-tiered minimums & High & SEC filing \cite{franklin} \\
Tokenized treasuries \$14.5--16bn & Medium & Trackers disagree; range quoted \\
Binance USYC July 2025, BUIDL Nov 2025 & Medium & One source dates USYC to 2024 \\
Deribit 2\% haircut on BUIDL & Medium & Single secondary source \\
Deribit/Crypto.com integration date & \textbf{Low} & Sources give Q1 2026 and July \\
Circle reserve \textasciitilde\$80--100m & \textbf{Low} & 2024 figures, may be stale \\
CFTC guidance 24 Sept 2026 & \textbf{Low} & \textbf{Not confirmed against primary source} \\
\bottomrule
\end{tabular}
\end{center}

\note{Items marked Low must be verified against primary sources — venue collateral documentation,
Circle's current documentation, and the CFTC's own releases — before this section is considered
final.}

\begin{thebibliography}{99}
\footnotesize\sloppy\emergencystretch=3em
\bibitem{circle} Circle, ``Circle Announces USDC Smart Contract for Transfers by BlackRock,''
Businesswire, 11 April 2024. \url{https://www.businesswire.com/news/home/20240411966052/en/Circle-Announces-USDC-Smart-Contract-for-Transfers-by-BlackRock}
\bibitem{binance} Cointelegraph, ``Binance adds BlackRock's BUIDL as off-exchange collateral for
institutional traders.'' \url{https://cointelegraph.com/news/binance-blackrock-buidl-off-exchange-collateral-traders}
\bibitem{deribit} Cointelegraph, ``Deribit, Crypto.com integrate BlackRock's BUIDL as trading
collateral.'' \url{https://cointelegraph.com/news/deribit-crypto-com-accepting-blackrock-buidl-collateral}
\bibitem{protofire} Protofire, ``Lending Against Tokenized Assets: Aave vs Morpho.''
\url{https://protofire.io/guides/lending-against-tokenized-assets/}
\bibitem{allium} Allium, ``Tokenized Treasuries Market Size: \$16B Across 91 Funds,'' 12 October 2026.
\url{https://www.allium.so/blog/tokenized-treasuries-market-size-16b-and-counting/}
\bibitem{cryptonomist} Cryptonomist, ``Tokenized US Treasuries Reach \$15.42B Value in 2026,''
18 October 2026. \url{https://en.cryptonomist.ch/2026/10/18/tokenized-us-treasuries-market-2026/}
\bibitem{utoday} U.Today, ``Cryptorank: Tokenized US Treasuries Hold \$14.7 Billion as RWA Market
Expands.'' \url{https://u.today/cryptorank-tokenized-us-treasuries-hold-147-billion-as-rwa-market-expands}
\bibitem{coincentral} CoinCentral, ``Circle USYC Overtakes BlackRock BUIDL as Largest Tokenized
Treasury Fund.'' \url{https://coincentral.com/circle-usyc-overtakes-blackrock-buidl-as-largest-tokenized-treasury-fund/}
\bibitem{ecokyc} Eco, ``Tokenized Treasury KYC Requirements 2026: BUIDL, OUSG, BENJI Access Tiers.''
\url{https://eco.com/support/en/articles/15254021-tokenized-treasury-kyc-requirements-2026-buidl-ousg-benji-access-tiers}
\bibitem{coinpaprika} Coinpaprika, ``BlackRock BUIDL Explained: Tokenized Fund by Securitize.''
\url{https://coinpaprika.com/education/blackrock-buidl-explained-tokenized-fund-by-securitize/}
\bibitem{franklin} Franklin Templeton Trust, Form 485BPOS, U.S.\ Securities and Exchange Commission.
\url{https://www.sec.gov/Archives/edgar/data/0001786958/000165558926000970/c485bpos.htm}
\bibitem{techflow} TechFlow, ``Understanding BlackRock's Tokenized Fund Experiment BUIDL: Model,
Current Status, and Challenges.'' \url{https://www.techflowpost.com/en-US/article/20976}
\bibitem{bitcoincom} Bitcoin.com News, ``Circle Unveils Smart Contract Service for Blackrock BUIDL
Holders to Convert to USDC.'' \url{https://news.bitcoin.com/circle-unveils-smart-contract-service-for-blackrock-buidl-holders-to-convert-to-usdc/}
\bibitem{rwaxyz} RWA.xyz, Tokenized Treasuries. \url{https://app.rwa.xyz/treasuries}
\bibitem{sherlock} Sherlock, ``Smart Contract Audit Pricing: A Market Reference for 2026.''
\url{https://www.sherlock.xyz/post/smart-contract-audit-pricing-a-market-reference-for-2026}
\bibitem{hlfees} Hyperliquid and dYdX published perpetual fee schedules, base tier: taker 4.5bp and
5.0bp respectively. \url{https://app.hyperliquid.xyz/docs}
\end{thebibliography}

\end{document}
